\clearpage
\section{The Polylogarithmic Three-Point Path Packet}

\subsection{Root-distribution compression}
Let $u^M=z$ and let $\zeta_M$ be a primitive $M$th root of unity. For $n\ge1$,
\[
 \boxed{
 \sum_{j=0}^{M-1}
 \operatorname{Li}_n(\zeta_M^ju)
 =
 M^{1-n}\operatorname{Li}_n(u^M).
 }
\]
Inside the disk of convergence, this follows by summing the defining series and using the root-of-unity filter. Analytic continuation extends the identity with compatible branches.

The inverse synthesis gain at weight $n$ is $M^{n-1}$. Through weights one, two, and three, the total gain is
\[
 \boxed{1\cdot M\cdot M^2=M^3.}
\]
The displayed coordinate rescaling has determinant $M^3$. It has not been identified here with an actual integral motivic or period-lattice map. That identification would require the missing arithmetic realization.

\subsection{Three-point functional closure}
Set
\[
 z=\frac ac,
 \qquad
 1-z=\frac bc,
 \qquad
 \beta=1-\frac1z=-\frac ba.
\]
The Rogers dilogarithm
\[
 L_2(z)
 =
 \operatorname{Li}_2(z)
 +
 \frac12\log z\log(1-z)
\]
satisfies
\[
 \boxed{
 L_2(z)+L_2(1-z)=\frac{\pi^2}{6}.
 }
\]
At weight three,
\[
 \boxed{
 \begin{aligned}
 &\operatorname{Li}_3(z)
 +\operatorname{Li}_3(1-z)
 +\operatorname{Li}_3(\beta)\\
 &\qquad=
 \zeta(3)
 +\frac16\log^3z
 +\frac{\pi^2}{6}\log z
 -\frac12\log^2z\log(1-z).
 \end{aligned}
 }
\]
For $0<z<1$, all three displayed arguments use their real branches. The dilogarithm complement identity proves the Rogers relation. The auxiliary identity
\[
 \Li_2(1-1/z)+\Li_2(1-z)=-\tfrac12\log^2z
\]
follows by differentiation and the limit at $z=1$. Using $\dd\Li_3(w)=\Li_2(w)\,\dd w/w$, the derivative of the displayed trilogarithm sum is
\[
 \frac{\pi^2}{6z}-\frac{\log z\log(1-z)}z+
 \frac{\log^2z}{2z(1-z)}.
\]
It is also the derivative of the right side. Both sides tend to $\zeta(3)$ as $z\to1$, proving the formula. The balanced Kummer coordinate is therefore exactly the third argument in the depth-three functional closure.

A fixed path packet is
\[
 \Pi_{\le3}(z)
 =
 \left(
 \log z,
 \log(1-z),
 L_2(z),
 \operatorname{Li}_3(z),
 \operatorname{Li}_3(1-z)
 \right).
\]
Its coordinate count is independent of $M$.

\subsection{Archimedean envelope}
For $0<z<1$,
\[
 |\log z|,
 |\log(1-z)|
 \le\log c.
\]
The positive dilogarithm and trilogarithm values on $(0,1)$ are bounded by $\zeta(2)$ and $\zeta(3)$. The three-point identity gives
\[
 \left|\operatorname{Li}_3\left(-\frac ba\right)\right|
 \le
 3\zeta(3)
 +
 \frac{\pi^2}{6}\log c
 +
 \frac23(\log c)^3.
\]
Thus the complex packet has only cubic logarithmic growth. The surrounding motivic and regulator language is informed by the polylogarithmic complexes of Goncharov, the regulator framework discussed by Zagier, and Borel's rank theorem~\cite{Goncharov,ZagierPolylog,Borel}.

This is a genuine nonlinear compression of the displayed complex data. It is not yet an integral compression. A fixed-rank integer lattice carrying a large torsion quotient must pay through covolume, denominator, algebraic degree, or a short vector.

\subsection{Fixed-rank integral uncertainty}
Let $A\in M_d(\Z)$ be nonsingular and suppose $D\mid\det A$. If every entry satisfies $|A_{ij}|\le B$, then Hadamard's inequality gives
\[
 D\le|\det A|\le d^{d/2}B^d.
\]
Hence
\[
 \boxed{
 B\ge\frac{D^{1/d}}{\sqrt d}.
 }
\]
A small complex period packet does not imply a small integral realization.

\subsection{Exact finite-place notch}
At an output prime $p^e\parallel c$,
\[
 \beta=1-\frac ca.
\]
Since $p\nmid a$,
\[
 v_p(\beta-1)=e.
\]
For odd $p$,
\[
 \boxed{
 v_p(\log_p\beta)=e.
 }
\]
The corresponding weighted dyadic identity holds on the principal branch when $e\ge2$.

Writing
\[
 \beta=-\prod_{q\mid ab}q^{\nu_q},
\]
we obtain the cross-prime logarithmic filter relation
\[
 \boxed{
 \log_p\beta
 =
 \sum_{q\mid ab}
 \nu_q\log_p\langle q\rangle_p.
 }
\]
The same depth-three path therefore has a small complex envelope and exact deep finite notches.

\subsection{Weight-three motivation and limitation}
Over $\Q$, Borel's rank theorem gives $\operatorname{rank}K_3(\Q)=0$ and $\operatorname{rank}K_5(\Q)=1$, the rational weights two and three. This makes weight three a plausible rational regulator layer. It does not assert the absence of all weight-two periods or of regulator directions over arbitrary number fields. This motivates the depth-three path packet. It does not prove that weight three suffices uniformly. A complete proof may require adaptive higher weights or a different motivic quotient.

\subsection{Conditional motivic transfer}
Let $K_\beta$ be the cyclic group generated by the balanced Kummer class. Suppose that for every primitive positive triple there exists an integral finite object $\mathcal H_3(a,b,c)$ and a canonical map
\[
 \Theta_3:K_\beta\longrightarrow\mathcal H_3(a,b,c)
\]
such that, for every $\eps>0$, constants independent of the triple give
\[
 |\ker\Theta_3|\le C_\eps R^\eps
\]
and
\[
 |\mathcal H_3(a,b,c)|\le C'_\eps R_aR_bR^\eps.
\]
Then the strong $abc$ conjecture follows after absorbing the logarithmic Kummer kernel.

The missing object must combine finite Kummer depth, the three-point additive relation, an integral motivic or syntomic lattice, and an archimedean period norm.

